% 《经济研究》(Economic Research Journal) LaTeX 模板
% 适配 GB/T 7714-2015 参考文献格式
% 格式：16开，双栏，中文，支持JEL分类号
%
% 编译方式：xelatex -> bibtex -> xelatex -> xelatex
%
% 依赖宏包（自动加载）：
%   ctex, xeCJK, natbib, amsmath, geometry, fancyhdr, setspace, graphicx, booktabs

% ─── 文档类 ────────────────────────────────────────────────────────────────
\documentclass[10.5pt, UTF8, twocolumn, landscape]{ctexart}
% 10.5pt: 经济研究正文字号
% twocolumn, landscape + geometry: 实现16开双栏排版

% ─── 页面布局 ─────────────────────────────────────────────────────────────
\usepackage[
    paperwidth=235mm,
    paperheight=165mm,
    top=20mm,
    bottom=20mm,
    left=15mm,
    right=15mm,
    headheight=10mm,
    footskip=8mm,
]{geometry}

% ─── 语言与字体 ───────────────────────────────────────────────────────────
\usepackage{xeCJK}
% 设置中文主体字体（系统须安装相应字体，fallback 由 ctex 自动处理）
%\setCJKmainfont{FandolSong-Regular.otf}[BoldFont={FandolSong-Bold.otf}]

% ─── 数学与符号 ────────────────────────────────────────────────────────────
\usepackage{amsmath, amssymb}
\usepackage{mathtools}   % enhanced amsmath
\usepackage{booktabs}    % professional tables
\usepackage{longtable}   % 长表格
\usepackage{graphicx}    % 插图
\graphicspath{{./figures/}}

% ─── 间距控制 ─────────────────────────────────────────────────────────────
\usepackage{setspace}
%\setstretch{1.5}          % 正文1.5倍行距（经济研究要求）
\setlength{\parindent}{2em}  % 首行缩进2字符
\setlength{\parskip}{0pt}    % 段间距

% ─── 页面样式 ─────────────────────────────────────────────────────────────
\usepackage{fancyhdr}
\fancypagestyle{JJYX}{
    \fancyhf{}
    \fancyhead[C]{\thepage}
    \fancyfoot[C]{\thepage}
    \renewcommand{\headrulewidth}{0pt}
    \renewcommand{\footrulewidth}{0pt}
}
\pagestyle{JJYX}

% ─── 超链接 ───────────────────────────────────────────────────────────────
\usepackage[colorlinks=true, linkcolor=black, citecolor=black, urlcolor=black]{hyperref}

% ─── 参考文献（GB/T 7714-2015 numeric）──────────────────────────────────────
\usepackage[numbers,sort&compress]{natbib}
% 样式说明：
%   gbt7714-2005.bst 或 gbt7714-plain.bst 配合 natbib 使用
%   如模板目录有 gbt7714-2005.bst，改用：
%   \bibliographystyle{gbt7714-2005}
\bibliographystyle{gbt7714-plain}

% ─── 标题宏定义 ──────────────────────────────────────────────────────────
% 经济研究要求：标题/作者/单位/摘要/关键词/中图分类号/JEL分类号
\makeatletter
% 作者、标题、关键词、摘要、分类号、JEL分类号存储
\def\@JJYXtitle#1{\gdef\@JJYX@title{#1}}\def\@JJYX@title{}
\def\@JJYXauthor#1{\gdef\@JJYX@author{#1}}\def\@JJYX@author{}
\def\@JJYXaffiliation#1{\gdef\@JJYX@affiliation{#1}}\def\@JJYX@affiliation{}
\def\@JJYXabstract#1{\gdef\@JJYX@abstract{#1}}\def\@JJYX@abstract{}
\def\@JJYXkeywords#1{\gdef\@JJYX@keywords{#1}}\def\@JJYX@keywords{}
\def\@JJYXclscode#1{\gdef\@JJYX@clscode{#1}}\def\@JJYX@clscode{}
\def\@JJYXjel#1{\gdef\@JJYX@jel{#1}}\def\@JJYX@jel{}

% 环境：摘要
\newenvironment{jjyxabstract}{
    \wuhao\sffamily\global\let\@JJYXabstract\@empty
    \par\noindent\textbf{摘要\quad}\wuhao
}{\par\vspace{\baselineskip}}

% 环境：关键词
\newenvironment{jjyxkeywords}{
    \par\noindent\textbf{关键词\quad}\wuhao
}{\par\vspace{\baselineskip}}

% 环境：分类号
\newenvironment{jjyxclscode}{
    \par\noindent\textbf{中图分类号\quad}\wuhao
}{\par\vspace{0.5\baselineskip}}

% 环境：JEL分类号
\newenvironment{jjyxjel}{
    \par\noindent\textbf{JEL分类号\quad}\wuhao
}{\par\vspace{\baselineskip}}

\makeatother

% 小五号字体命令
\newcommand{\wuhao}{\fontsize{9pt}{13pt}\selectfont}
% 标题区字号
\newcommand{\JJYXtitlefont}{\fontsize{16pt}{22pt}\bfseries\selectfont}
% 作者区字号
\newcommand{\JJYXauthorfont}{\fontsize{10.5pt}{16pt}\selectfont}

% ─── 标题页环境 ──────────────────────────────────────────────────────────
% 使用方法：
%   \begin{jjYXtitle}{标题}{作者}{单位}
%   \begin{jjYXabstract} 摘要内容 \end{jjYXabstract}
%   \begin{jjYXkeywords} 关键词1；关键词2；关键词3 \end{jjYXkeywords}
%   \begin{jjYXclscode} F0；F03 \end{jjYXclscode}
%   \begin{jjYXjel} JEL: G00; O40 \end{jjYXjel}
%   \end{jjYXtitle}

\NewDocumentEnvironment{jjYXtitle}{mmm}{
    % 参数：{标题}{作者信息}{单位信息}
    \twocolumn[%
        \centering
        % 标题
        \JJYXtitlefont\@JJYX@title\par
        \vspace{6pt}
        % 作者
        {\JJYXauthorfont #2}\par
        \vspace{3pt}
        % 单位
        {\JJYXauthorfont #3}\par
        \vspace{6pt}
        % 摘要栏（单栏排版）
        \hspace{2em}\begin{minipage}{0.9\linewidth}
            % 摘要
            \wuhao\textbf{摘要}\quad\@JJYX@abstract\par
            \vspace{3pt}
            % 关键词
            \wuhao\textbf{关键词}\quad\@JJYX@keywords\par
            % 中图分类号
            \ifx\@JJYX@clscode\@empty\else
                \wuhao\textbf{中图分类号}\quad\@JJYX@clscode\par
            \fi
            % JEL分类号
            \ifx\@JJYX@jel\@empty\else
                \wuhao\textbf{JEL分类号}\quad\@JJYX@jel\par
            \fi
        \end{minipage}
        \vspace{8pt}
    ]%
}{}

% ─── 标题快捷命令（单命令方式）────────────────────────────────────────────
% 使用方法：
%   \JJYXsettitle{标题}
%   \JJYXsetauthor{作者}   % 盲审时注释掉
%   \JJYXsetaffiliation{单位}
%   \begin{jjYXabstract} 摘要 \end{jjYXabstract}
%   \begin{jjYXkeywords} 关键词1；关键词2 \end{jjYXkeywords}

\NewDocumentCommand{\JJYXsettitle}{m}{\@JJYXtitle{#1}}
\NewDocumentCommand{\JJYXsetauthor}{m}{\@JJYXauthor{#1}}
\NewDocumentCommand{\JJYXsetaffiliation}{m}{\@JJYXaffiliation{#1}}
\NewDocumentCommand{\JJYXsetabstract}{m}{\@JJYXabstract{#1}}
\NewDocumentCommand{\JJYXsetkeywords}{m}{\@JJYXkeywords{#1}}
\NewDocumentCommand{\JJYXsetclscode}{m}{\@JJYXclscode{#1}}
\NewDocumentCommand{\JJYXsetjel}{m}{\@JJYXjel{#1}}

% ─── 正文环境 ─────────────────────────────────────────────────────────────
% 一级标题：宋体加粗
\ctexset{section/format={\centering\zihao{-3}\CJKfamily{zhsong}\bfseries}}
% 二级标题：宋体
\ctexset{subsection/format={\zihao{4}\CJKfamily{zhsong}}}
% 三级标题：宋体
\ctexset{subsubsection/format={\zihao{-4}\CJKfamily{zhsong}}}

% ─── 图表浮动体 ────────────────────────────────────────────────────────────
\usepackage{flafter}    % 图表放在引用之后
\usepackage{booktabs}   % 三线表
\usepackage{sidecap}    % 图注在侧
% 表格字号
\DeclareFontFamily{OT1}{fakesc}{}
\usepackage{caption}
\captionsetup{font=small, skip=6pt}

% ─── 定理/证明环境 ─────────────────────────────────────────────────────────
\newtheorem{hypothesis}{假设}
\newtheorem{assumption}{假定}
\newtheorem{definition}{定义}
\newtheorem{prop}{性质}
\newtheorem*{proof}{证明}

% ─── 附录格式 ─────────────────────────────────────────────────────────────
\ctexset{section/aftertitle={\par}}{}

% ─────────────────────────────────────────────────────────────────────────
\begin{document}

% ══════════════════════════════════════════════════════════════════════════════
% 请填写以下信息（盲审时请删除作者和单位信息）
% ══════════════════════════════════════════════════════════════════════════════

\JJYXsettitle{论文标题（中文，不超过20字）}

% 盲审时注释掉以下两行
\JJYXsetauthor{作者一$^{1}$  ~~作者二$^{2}$ ~~作者三$^{1}$}
\JJYXsetaffiliation{$^{1}$北京大学光华管理学院  $^{2}$清华大学五道口金融学院}

% 摘要（不超过300字）
\begin{jjYXabstract}
本文基于XX数据，采用XX方法，实证研究了XX问题。研究发现：（1）……；（2）……；（3）……。这一发现对理解……具有重要意义，并对我……政策的制定提供了经验证据。
\end{jjYXabstract}

% 关键词（3-5个）
\begin{jjYXkeywords}
关键词一；关键词二；关键词三；关键词四；关键词五
\end{jjYXkeywords}

% 中图分类号（http://ztflh.xhma.com/）
\begin{jjYXclscode}
F0；F03
\end{jjYXclscode}

% JEL分类号（https://www.aeaweb.org/jel/jel-class-series）
\begin{jjYXjel}
JEL: G00; O40; C50
\end{jjYXjel}

% ══════════════════════════════════════════════════════════════════════════════
% 正文
% ══════════════════════════════════════════════════════════════════════════════

\section{引言}
\label{sec:intro}

\section{文献综述与研究假设}
\label{sec:lit}

\subsection{文献综述}

\subsection{研究假设}
基于以上分析，本文提出如下待检验的研究假设：

\begin{hypothesis}
\label{hyp:h1}
\textbf{假设1：}……（说明假设内容）。
\end{hypothesis}

\begin{hypothesis}
\label{hyp:h2}
\textbf{假设2：}……（说明假设内容）。
\end{hypothesis}

\section{研究设计}
\label{sec:design}

\subsection{样本与数据}
本文以……为研究样本，时间跨度为……至……。数据来源于……。

\subsection{变量定义}

\subsubsection{被解释变量}
$Y_{it}$ 表示……，数据来源于……。

\subsubsection{解释变量}
核心解释变量$X_{it}$ 为……。

\subsubsection{控制变量}
参考已有文献（\cite{xxx}），本文还控制了以下变量……。

\begin{table}[htbp]
\centering
\caption{变量定义表}
\label{tab:var}
\wuhao
\begin{tabular}{p{2.2cm}p{2.5cm}p{5.5cm}}
\toprule
变量类型 & 变量名称 & 变量定义 \\
\midrule
因变量 & $Y$ & ……（数据来源：XX数据库）\\
\midrule
自变量 & $X$ & ……（指标构造方法）\\
\midrule
控制变量 & $Z_1$ & ……\\
& $Z_2$ & ……\\
& $Year$ & 年份虚拟变量\\
\midrule
工具变量 & $IV$ & ……（工具变量选择依据）\\
\bottomrule
\end{tabular}
\end{table}

\subsection{模型设定}
为检验假设1，本文构建如下双向固定效应回归模型：

\begin{equation}
\label{eq:baseline}
Y_{it} = \alpha + \beta X_{it} + \gamma Z_{it} + \mu_i + \lambda_t + \varepsilon_{it}
\end{equation}

其中，$\mu_i$ 和 $\lambda_t$ 分别为企业和年份固定效应，$\varepsilon_{it}$ 为误差项。

\section{实证结果}
\label{sec:results}

\subsection{描述性统计}

\begin{longtable}{ccccccc}
\caption{描述性统计}
\label{tab:desc}
\\\toprule
变量 & 观测数 & 均值 & 标准差 & 最小值 & 中位数 & 最大值 \\
\midrule
\endfirsthead
\multicolumn{7}{c}{\textbf{续表\thetable\ 描述性统计}}\\
\toprule
变量 & 观测数 & 均值 & 标准差 & 最小值 & 中位数 & 最大值 \\
\midrule
\endhead
\bottomrule
\endfoot
$Y$ & 1000 & 0.05 & 0.20 & -0.50 & 0.03 & 0.60 \\
$X$ & 1000 & 0.30 & 0.15 & 0.10 & 0.28 & 0.70 \\
$Z_1$ & 1000 & 5.20 & 1.50 & 1.00 & 5.00 & 9.00 \\
\bottomrule
\end{longtable}

\subsection{基准回归分析}
表~\ref{tab:reg} 报告了基准回归结果。

\begin{table}[htbp]
\centering
\caption{基准回归结果}
\label{tab:reg}
\wuhao
\begin{tabular}{lccc}
\toprule
 & (1) & (2) & (3) \\
\midrule
$X$ & 0.05** & 0.06*** & 0.04** \\
    & (0.02) & (0.02) & (0.02) \\
$Z_1$ & & 0.10*** & 0.08** \\
      &      & (0.03)  & (0.03) \\
\midrule
常数项 & 0.01 & 0.02 & 0.03 \\
       & (0.02) & (0.02) & (0.02) \\
\midrule
控制变量 & 否 & 是 & 是 \\
企业固定效应 & 否 & 否 & 是 \\
年份固定效应 & 否 & 否 & 是 \\
\midrule
观测数 & 1000 & 1000 & 1000 \\
$R^2$ & 0.15 & 0.18 & 0.22 \\
\bottomrule
\multicolumn{4}{l}{\small 注：*** p<0.01, ** p<0.05, * p<0.1。括号内为聚类稳健标准误（聚类到企业）。}
\end{tabular}
\end{table}

\subsection{内生性处理}
为缓解遗漏变量偏误和反向因果导致的内生性问题，本文采用……方法进行检验。

\subsection{稳健性检验}

\subsubsection{替换核心变量}
为保证结果的稳健性，本文采用……替换核心解释变量。

\subsubsection{改变样本范围}
剔除……样本后，主要结论依然成立。

\section{进一步分析}
\label{sec:进一步}

\section{结论与启示}
\label{sec:concl}

本文基于……数据，实证研究了……问题。研究发现：（1）……；（2）……；（3）……。

本文的研究启示在于：第一，……；第二，……；第三，……。未来研究可以从……角度进一步深化。

\section*{参考文献}
\addcontentsline{toc}{section}{参考文献}
\bibliography{references}

% ══════════════════════════════════════════════════════════════════════════════
% 附录（如有）
% ══════════════════════════════════════════════════════════════════════════════
\clearpage
\appendix
\section*{附录}
\addcontentsline{toc}{section}{附录}

\end{document}
